\chapter{Sviluppi futuri}
\label{ch:sviluppi_futuri}

Sulla scorta delle valutazioni euristiche e degli esperimenti con gli utenti descritti nei capitoli precedenti, sono state identificate alcune aree di miglioramento che potrebbero essere affrontate nelle prossime fasi di sviluppo del progetto \textit{Guided Tours}.

Di seguito si riportano i principali sviluppi futuri previsti:

\begin{itemize}

    \item \textbf{Gestione avanzata e personalizzazione per l'amministratore:} 
    Concedere all'amministratore un controllo più granulare del sistema, consentendo non solo di modificare i luoghi o le guide presenti, ma di personalizzare componenti grafiche o di notifica lato utente. Questo andrebbe a coprire alcuni casi d'uso limite evidenziati nelle ispezioni euristiche e a semplificare la customizzazione senza l'intervento diretto degli sviluppatori.
    
    \item \textbf{Ottimizzazione della disponibilità guide e volontari:}
    In riferimento ad alcune osservazioni degli utenti in fase di test, sarebbe interessante potenziare le interfacce dedicate all'inserimento delle disponibilità. Un possibile miglioramento vedrebbe l'introduzione di strumenti di selezione multipla o ricorrente sul calendario, permettendo al volontario di confermare la propria presenza per "tutti i fine settimana" con una singola interazione.
    
    \item \textbf{Sviluppo di un'app mobile dedicata:}
    Potrebbe risultare molto vantaggiosa l'introduzione di un'applicazione mobile nativa destinata ai visitatori. Questa app potrebbe offrire funzionalità interattive durante la visita guidata \textit{in loco}: lettura di QR Code, supporto ad audio-guide multilingua basate sulla posizione GPS, ricevimento di notifiche *push* per cambi d'orario, e recensioni istantanee.

    \item \textbf{Ampliamento dell'accessibilità e Internazionalizzazione:}
    Sebbene l'interfaccia sia ora progettata seguendo forti principi di contrasto visivo e leggibilità, i futuri sviluppi dovrebbero mirare all'adeguamento totale agli standard WCAG 2.1 AAA. Elementi di supporto come lo *screen reading* ottimizzato e un sistema robusto di traduzione istantanea renderebbero la piattaforma totalmente scalabile a livello turistico globale.
    
    \item \textbf{Miglioramento ed espansione del motore di ricerca (Faceted Search):}
    L'attuale filtro delle visite guidate ha aumentato drasticamente le performance di conversione nei test (Compito 2), ma potrebbe essere arricchito. Introdurre mappe interattive in cui evidenziare le zone ad alta densità di tour, o l'integrazione di tag semantici ("Kid-friendly", "Accessibile in carrozzina") ridurrebbe il carico di memoria di lavoro (HSA 3) nell'utente in fase di decisione.

\end{itemize}

\section{Sviluppi Futuri e Workaround per Problemi Euristici Irrisolti}
Durante la fase di refactoring, il team ha volutamente tralasciato la correzione di cinque problemi emersi nell'ispezione euristica, classificandoli temporaneamente come \textit{out-of-scope} a causa dell'elevata complessità implementativa in rapporto alle tempistiche. Di seguito si tracciano soluzioni metodologiche e \textit{workaround} di UX per tre di queste criticità, predisponendole per i futuri cicli di sviluppo iterativo:

\begin{itemize}
    \item \textbf{Mancanza di selezione multipla nel calendario:}
    \newline\textit{Problema identificato:} Nelle interfacce di dichiarazione disponibilità, la necessità di reiterare innumerevoli click disgiunti per profilare giorni contigui sovraccarica le procedure operative dal lato delle guide, aumentando stanchezza e potenziale errore umano in aperta violazione delle euristiche su flessibilità ed efficienza.
    \newline\textit{Sviluppo Futuro e Workaround:} L'evoluzione strutturale ideale prevedrà la trasposizione funzionale del widget attuale per l'adozione di una libreria DOM orientata alla selezione multipla temporale integrata in logiche \textit{batch-processing} (es. plugin \textit{Flatpickr} multiselect). Tramite la gestualità di \textit{click-and-drag}, il carico cognitivo si annullerebbe radunando la sequenza all'interno di un solo carico per l'HTTP POST. La soluzione \textit{workaround} provvisoria per le fasi critiche, consisterà nell'esposizione di due meri input (``Data di Inizio'' e ``Data di Fine'') con uno script di backend preposto alla deduzione logico-matematica progressiva delle date intermedie.

    \item \textbf{Assenza di notifica per annullamento visite:}
    \newline\textit{Problema identificato:} Nel configuratore aziendale, l'epurazione di uno slot di visita da parte di un Amministratore non scaturisce riscontri attivi sui dispositivi degli iscritti, mantenendo silente lo stato e violando nettamente l'euristica di visibilità e trasparenza dello stato del sistema utente.
    \newline\textit{Sviluppo Futuro e Workaround:} L'architettura correttiva vedrà implementato sul framework Laravel un robusto pattern logico \textit{Publisher/Subscriber}: l'intercettazione del task \textit{soft-delete} della risorsa imposterà un transito elaborativo all'interno di una \textit{Queue Job} che diramerà \textit{email transazionali d'allerta} verso ogni singolo fruitore iscritto alla seduta. Nell'immediato volendo rifuggire \textit{technical debt} si istruisce un solido \textit{alert block} temporaneo all'interno dei modali di approvazione dei \textit{delete}: \textit{<<Attenzione: Tale cancellazione è irreversibile ma per il momento non invia notifiche automatiche via mail. Si prega di farsi manualmente carico della profilazione d'invio avvisi agli associati allo slot anzitempo il task.>>}
\end{itemize}
